RECOMENDATIONS===================================================

Recommendation 1 Explore human centred and exteprise level constituents.
Most SoTS studies focus on physical assets represented by fully autonomous digital twins, with the majority of evaluated systems assuming minimal human intervention and optimization-driven autonomy. In contrast, cyber-physical-human systems are rarely examined in the existing literature, despite their reliance on human decision-making and involvement. Enterprise-level constituents that introduce organizational structure and policy constraints are also underexplored. Future work should broaden the range of constituent system types studied and explicitly examine alternative autonomy configurations, especially in safety-critical or socio-technical contexts where accountability and coordination across human and technical actors are essential. This underrepresentation has important implications for the scope of SoTS applicability. Many real-world systems of systems—particularly in public infrastructure, healthcare, governance, and enterprise settings—are shaped as much by human decision-making, institutional constraints, and organizational processes as by technical coordination. By centering primarily on physical assets and autonomous DTs, current SoTS research risks overlooking challenges related to accountability, policy compliance, and human-system interaction. The results therefore point to a structural gap between the systems most frequently studied in SoTS research and the socio-technical systems that motivate much of SoS theory.


Recommendation 2 (Intent-aware and decentralized SoTS design).
Across the reviewed studies, SoTS designs overwhelmingly emphasize centralized coordination and operational controllability. This is reflected in the dominance of Acknowledged and Directed architectural patterns, the prevalence of fully autonomous DTs, and the strong focus on optimization-driven motivations.
Such configurations enable global system awareness, coordinated decision-making, and efficient execution, which align well with domains like manufacturing, automotive, and infrastructure systems. However, this emphasis also shapes the kinds of SoS dimensions that are actively supported: distribution, independence, interdependence, and interoperability are frequently addressed, while reconfiguration, evolution, and emergence receive comparatively limited attention.

current SoTS research implicitly trades SoS dynamism for predictability and control. While this trade-off may be appropriate for optimization-centric or safety-critical applications, it constrains the applicability of SoTS to domains characterized by uncertainty, changing system boundaries, or evolving stakeholder goals.

Our results show a strong concentration of SoTS research in manufacturing contexts, where all SoTS classifications are well represented, while domains such as healthcare, logistics, cybersecurity, and enterprise systems remain  underexplored. Moreover, studies in these non-industrial domains are largely limited to directed or acknowledged SoTS configurations, with collaborative and virtual SoTS appearing only sporadically. Future research should intentionally target underrepresented domains to assess whether existing SoTS assumptions, architectures, and coordination mechanisms generalize beyond industrial automation settings.

The prevalence of combining DTs into a SoS, alongside the dominance of centralized Acknowledged and Directed architectures, suggests an implicit design bias toward bottom-up integration and centralized coordination. Future studies should explicitly justify the chosen SoTS intent and architectural configuration, and further explore decentralized Collaborative and Virtual SoTS to better understand autonomy, negotiation, and emergent behavior in highly distributed environments.

Although many systems are framed as SoS, emergent behavior, reconfiguration, and long-term evolution are only partially addressed or omitted entirely. Future research should explicitly analyze and evaluate these dynamic properties, for example through emergence-aware simulations, scenario-based validation, or longitudinal studies that capture system adaptation over time.



Recommendation 3 (Formalization and validation of reliability and secuirty).
While reliability is frequently cited as a motivation or architectural concern, it is rarely formalized or empirically validated. Reliability and security emerge as recurring concerns in SoTS research, yet they are addressed predominantly at the architectural level rather than through formal modeling or systematic evaluation. Future SoTS research should adopt explicit reliability models and validation techniques, such as fault injection, simulation-based evaluation, or formal analysis, to move beyond architectural intuition. . While SoTS inherently introduce additional sources of uncertainty—due to distribution, autonomy, and coordination across constituents—most studies stop short of making these uncertainties explicit or quantifying their impact on system behavior.

Our results show that SoTS evaluations are predominantly conducted through simulation and prototyping under nominal or expected operating conditions. While these validation approaches demonstrate functional feasibility and coordination capabilities, they rarely examine system behavior under stress, degradation, or abnormal operating scenarios. Even when reliability or security is mentioned, evaluation typically focuses on architectural mitigation rather than systematic exploration of failure modes, cascading effects, or degraded performance across interacting constituents.

Given the distributed and interdependent nature of SoTS, system behavior under non-nominal conditions is especially critical. Failures, delayed updates, partial observability, and inconsistent constituent behavior can propagate in ways that are not apparent during steady-state operation. Future SoTS research should therefore broaden evaluation strategies to explicitly include stress scenarios, such as partial system failures, communication disruptions, conflicting objectives, or degraded sensing. Evaluating SoTS under such conditions would provide deeper insight into their robustness and coordination mechanisms, and better reflect the challenges faced in real-world systems of systems.



Recommendation 4 (Advancing maturity through evaluation and standardization).
The dominance of low-to-mid TRL prototypes and validation-oriented studies indicates a maturity gap in SoTS research. Validation is typically conducted through simulation or prototyping, with relatively few studies assessing systems in operational or organizational contexts. The results therefore suggest that the maturation of SoTS research depends not only on proposing new architectures or frameworks, but also on grounding these proposals in empirical evidence that captures the complexity and variability of real systems of systems. Future work should prioritize evaluation research, including industrial case studies and align with industery standards. evelop architectural specifications and reference
implementations for SoTS to ease their engineer-
ing and to allow higher levels of maturity in their
research and development.


Recommendation 5 (Improve reporting consistency and comparability across SoTS studies)

Our results reveal substantial heterogeneity in how SoTS studies report system characteristics, design decisions, and evaluation scope. Key aspects such as autonomy configuration, supported DT services, SoS dimensions, treatment of emergence, and non-functional properties are often described unevenly or only implicitly. In many cases, these characteristics must be inferred rather than explicitly stated, which complicates cross-study comparison and synthesis. This variability is evident across multiple RQs, where similar systems differ widely in the depth and clarity of reported information.

This lack of reporting consistency limits the cumulative value of SoTS research, as it hinders systematic comparison of architectures, design choices, and outcomes across domains. Future work would benefit from more structured and explicit reporting of core SoTS characteristics, including constituent types, architectural intent, autonomy levels, DT services, SoS dimensions, and evaluation context. Establishing clearer reporting conventions would not constrain design freedom, but would significantly improve interpretability, reproducibility, and the ability to build evidence-based insights across the growing body of SoTS literature.



Recommendation 7 (Critically examine domain-specific assumptions and transferability)

Our results show a strong concentration of SoTS research within a small set of application domains, particularly manufacturing, automotive, and industrial infrastructure. In these domains, assumptions such as stable system structure, well-defined operational goals, and centralized coordination are often valid and reinforce each other. As a result, SoTS architectures, DT services, and autonomy configurations observed in the literature are frequently shaped by domain-specific constraints rather than by generalizable SoTS principles. This concentration risks conflating domain suitability with architectural optimality.

Consequently, it remains unclear to what extent prevailing SoTS design patterns transfer to domains with fundamentally different characteristics, such as healthcare, public administration, cybersecurity, or logistics, where objectives may be conflicting, boundaries less stable, and human or organizational factors more prominent. Future SoTS research should therefore explicitly assess the domain-dependence of proposed approaches and investigate how architectural assumptions change when applied outside industrial automation contexts. Making domain assumptions explicit would help distinguish between SoTS patterns that are broadly applicable and those that are effective primarily within a narrow set of domains.



\begin{recoframe}{Recommendation Template}
Prioritize evaluation research, including comparative and longitudinal studies, and explicit standard alignment in SoTS to advance maturity, improve interoperability, and support technology transfer.
\end{recoframe}








% =========================================================================================

Recommendation 1 — Adopt more flexible and decentralized SoTS design
Across the reviewed studies, SoTS designs emphasize relatively centralized coordination, as reflected in the dominance of Acknowledged and Directed architectures. While these configurations support efficient task execution, they place limited emphasis on SoS principles associated with dynamic interaction and runtime adaptation. As a result, little attention is given to reconfiguration, evolution, and emergence. Future work should intentionally target underrepresented SoS properties, and further explore Collaborative and Virtual SoTS to better understand decentralization, negotiation, and emergent behavior in dynamic and uncertain environments.


Recommendation 2 — Evaluate SoTS under non-ideal operating conditions
While reliability and security are frequently cited as motivations or architectural concerns in SoTS research, they are rarely formalized or empirically evaluated. Most studies assess these properties through architectural reasoning or simulation-based validation under idealized conditions, with limited examination of failure propagation, partial observability, or degraded constituent behavior. Given the distributed and interdependent nature of SoTS, such non-ideal operating conditions can interact in non-obvious ways and significantly influence system behaviour. Future SoTS research should therefore incorporate systematic evaluation under failure, delay, uncertainty, and inconsistency to better understand how SoTS behave beyond normal operation.

Recommendation 3 — Examine domain dependence and cross-domain transferability of SoTS designs (too similar to Rec 4, ask for advice, maybe combine)
SoTS research is strongly concentrated in manufacturing, automotive, and industrial infrastructure domains. As a result, commonly observed SoTS architectures, coordination strategies, and autonomy assumptions may reflect domain-specific constraints rather than generally applicable design principles. It remains unclear to what extent these patterns transfer to domains characterized by conflicting goals, evolving boundaries, regulatory oversight, or strong human and organizational influences. Future research should explicitly assess domain dependence and investigate how SoTS assumptions and architectural choices change when applied outside industrial automation contexts.

Recommendation 4 — Broaden SoTS constituent models to include human-centred and enterprise-level systems
Most SoTS studies focus on physical assets represented by fully autonomous digital twins, with minimal human involvement. In contrast, cyber-physical-human systems and enterprise-level constituents, where system behaviour is shaped by human decision-making, are rarely examined. This narrow constituent focus limits the applicability of SoTS to many real-world systems of systems, particularly in domains such as urban infrastructure, healthcare, and public-sector operations. Future SoTS research should broaden the range of constituent types studied and explicitly investigate alternative autonomy configurations, especially in socio-technical and safety-critical contexts where accountability, coordination, and human oversight are essential.

Recommendation 5 — Advance SoTS maturity through stronger evaluation research and shared architectures (Crosses Rec 1 and 2)
The predominance of low-to-mid TRL prototypes and validation-oriented studies indicates a maturity gap in SoTS research. Evaluation is most often conducted through simulations or short-lived prototypes, with relatively few studies examining sustained operation, real-world constraints, or long-term system evolution. Advancing SoTS maturity requires increased emphasis on case study evaluation alongside stronger alignment with relevant standards. In particular, developing shared architectural specifications, reference models, and reusable components would support reproducibility, enable cross-study comparison, and improve interoperability across heterogeneous and composed SoTS systems.

% ==========================================================

